\section{Related Work}

Several prior works have investigated the use of LLMs for decision-making. In one category, numerous studies have deployed LLMs directly as agents in decision-making problems such as games~\citep{yao2023react,brooks2024large,shinn2024reflexion,wang2023voyager,xi2023rise}. However, fewer works have systematically evaluated LLMs' capabilities in general decision-making setup, especially in relation to classical concepts in decision-making like exploration. Our work extends the research of \citet{krishnamurthy2024can}, who examined LLMs' exploration capabilities in small-scale MAB tasks. Their findings, which showed positive results only with substantial intervention, are consistent with our broader analysis across both MAB and CB tasks at various scales. \citet{mirchandani2023large,rahn2024controlling,felicioni2024importance} also evaluated the ability of LLMs to learn in-context and solve bandit-like decision-making problems.

Another relevant line of research focuses on in-context learning for decision-making and reinforcement learning (RL) with domain-specific transformers. \citet{laskin2022context} distilled demonstrations from RL algorithms into a transformer and showed that it learns to imitate the RL process to solve new RL tasks. Similarly, \citet{lee2024supervised} trained transformers with optimal action labels, showing that the model learns to execute posterior sampling for RL~\citep{osband2013more} in-context. This area has been further studied by~\citet{raparthy2023generalization,lin2023transformers,bai2023transformers}. However, these studies focus on domain-specific decision-making, whereas our paper examines general-purpose decision-making capabilities in language models. 
Our inference-time algorithm-guided support shares a similar conceptual framework with recent efforts to align LLMs at inference time. These include employing explicit value functions as prefix scorers~\citep{mudgal2023controlled}, and leveraging both implicit and explicit value functions to guide decoding~\citep{liu2024inference}. 
In the realm of algorithm distillation, much of the research on LLMs has concentrated on chain-of-thought (CoT) reasoning~\citep{wang2023scott, li2023symbolic}, while \cite{gandhi2024stream} focused on search and backtracking. Our work focuses on distilling a more complex class of algorithms that involve uncertainty estimation and linear regression.

Our inference-time algorithm-guided support shares a similar conceptual framework with recent efforts to align LLMs at inference time. These include employing explicit value functions as prefix scorers that trained via KL-regularized RL~\citep{mudgal2023controlled}, and leveraging both implicit and explicit value functions to guide decoding at the token and chunk levels at inference time~\citep{liu2024inference}. 
In the realm of algorithm distillation, much of the research on LLMs has concentrated on chain-of-thought (CoT) reasoning~\citep{wang2023scott, li2023symbolic}, while \cite{gandhi2024stream} focused on search and backtracking. Most methods involve distilling outputs from a "teacher" model—either a larger model or a slower, system-2 variant of the same model that employs various inference-time techniques, such as search and self-consistency—into a student model~\citep{yu2024distilling}. Instead, our approach leverages diverse optimal trajectories directly from classical algorithms, allowing for the efficient generation of abundant training data.

The motivation of our work is also inspired by the cognitive science literature on human. Prior studies have found that humans balance exploration and exploitation through a mix of directed and random exploration~\citep{wilson2014humans}, and that these behaviors are supported by distinct neural mechanisms~\citep{daw2006cortical}. \citet{gershman2018deconstructing} further decomposes human exploration into algorithmic components. These works provide future directions for evaluating LLM's behavior in addition to reward~\cite{pan2025large}. Additionally, efficient exploration also benefits other areas of RL, such as multi-agent collboration~\citep{qu2024choices} and preference optimization~\citep{bai2025online}.

